\SourcePage{383}{388}
\section{Positronium}

The formulas obtained in the preceding paragraph can be applied to positronium---a hydrogen-like system of electron and positron.

In the centre-of-inertia system the operators of momenta of the electron and positron: $\hat{\mathbf{p}}_- = -\hat{\mathbf{p}}_+ \equiv \hat{\mathbf{p}}$, where $\hat{\mathbf{p}} = -i\hbar\boldsymbol{\nabla}$ is the operator of the momentum of relative motion corresponding to the relative radius vector $\mathbf{r} = \mathbf{r}_- - \mathbf{r}_+$. The full Hamilto-

\SourcePage{384}{389}
nian of positronium\,\footnote{In ordinary units.}
\begin{equation}
\begin{split}
\hat{H} &= \frac{\hat{\mathbf{p}}^2}{m} - \frac{e^2}{r} + \hat{V}_1 + \hat{V}_2 + \hat{V}_3,\\
\hat{V}_1 &= \frac{\hat{\mathbf{p}}^4}{4m^3c^2} + 4\pi\mu_0^2\delta(\mathbf{r}) - \frac{e^2}{2m^2c^2r}\left\{\hat{\mathbf{p}}^2+\frac{\mathbf{r}(\mathbf{r}\hat{\mathbf{p}})\hat{\mathbf{p}}}{r^2}\right\},\\
\hat{V}_2 &= 6\mu_0^2\frac{1}{r^3}\hat{\mathbf{l}}\hat{\mathbf{S}},\\
\hat{V}_3 &= 6\mu_0^2\frac{1}{r^3}\left\{\frac{(\hat{\mathbf{S}}\mathbf{r})(\hat{\mathbf{S}}\mathbf{r})}{r^2}-\frac{1}{3}\hat{\mathbf{S}}^2\right\}+4\pi\mu_0^2\left(\frac{7}{3}\hat{\mathbf{S}}^2-2\right)\delta(\mathbf{r}).
\end{split}
\tag{84.1}
\end{equation}
Here $\mu_0 = e\hbar/(2mc)$ is the Bohr magneton, $\hat{\mathbf{l}} = [\mathbf{r}\hat{\mathbf{p}}]$ is the operator of orbital angular momentum, $\hat{\mathbf{S}} = (\boldsymbol{\sigma}_+ + \boldsymbol{\sigma}_-)/2$ is the operator of total spin of the system (its square $\hat{\mathbf{S}}^2 = \frac{1}{2}(3 + \boldsymbol{\sigma}_+\boldsymbol{\sigma}_-)$). In $\hat{V}_1$ are included all correction terms of purely orbital character; $\hat{V}_2$ is the spin-orbit interaction; $\hat{V}_3$ includes the spin-spin and ``annihilation'' interactions.

The ``unperturbed'' Hamiltonian
\[
\hat{H} = \frac{\hat{\mathbf{p}}^2}{m} - \frac{e^2}{r}
\]
differs, naturally, from the Hamiltonian of the hydrogen atom only by the replacement of the electron mass by the reduced mass $m/2$. The energy levels of positronium are therefore half (in absolute value) of the levels of the hydrogen atom:
\begin{equation}
E = -\frac{me^4}{4\hbar^2 n^2}
\tag{84.2}
\end{equation}
($n$ is the principal quantum number).

The remaining terms in~(84.1) lead to a splitting of the levels~(84.2)---the appearance of fine structure. The arising levels are classified first of all by the values of the total angular momentum $j$. We also see that the spin operators of the particles enter in the Hamiltonian~(84.1) only in the form of the sum~$\hat{\mathbf{S}}$. This means that the Hamiltonian commutes with the operator of the square of the total spin $\hat{\mathbf{S}}^2$, i.e.\ the total spin continues to be conserved also in the approximation under consideration (second in $1/c$). Therefore the energy levels of positronium can also be classified by the total spin, taking the values $S = 0$ and $S = 1$. The levels with spin~0 are called levels of \textit{parapositronium}, and the levels with spin~1---levels of \textit{orthopositronium}.

Let us emphasise that the conservation of total spin in positronium is in fact an exact law, not connected with this or that approximation by $1/c$; it follows from the

\SourcePage{385}{390}
$CP$-invariance of electromagnetic interactions. Positronium represents a truly neutral system, and therefore its states are characterised by definite charge and combined parities. The latter is equal to $(-1)^{S+1}$ (see Problem to~\S\,27); since $S$ can take only the two values, 0 and~1, the conservation of combined parity is equivalent to conservation of total spin.

For $S = 0$ the total moment $j$ coincides with the orbital. For spin $S = 1$ and given $j$ the number $l$ takes the values $j$, $j\pm 1$, correspondingly each level ($n$, $j$) of orthopositronium splits, generally speaking, into three levels. Since the two values $l = j$ and $l = j\pm 1$ correspond to different parities, the Hamiltonian has no matrix elements connecting these states. However the operator of perturbation (the first member in $\hat{V}_3$), generally speaking, has non-diagonal matrix elements connecting states with $l = j + 1$ and $l = j - 1$, so that when this number $l$ loses, of course, its strict meaning of the orbital angular momentum.

Specific properties are possessed by the Zeeman effect in positronium (\textit{B.\,B.\,Berestetskii, I.\,Ya.\,Pomeranchuk}, 1949). The orbital magnetic moment of positronium is always zero: since in positronium $[\mathbf{r}_+\hat{\mathbf{p}}_+] = [\mathbf{r}_-\hat{\mathbf{p}}_-]$, the operator
\[
\hat{\boldsymbol{\mu}}_l = \mu_0([\mathbf{r}_+\hat{\mathbf{p}}_+] - [\mathbf{r}_-\hat{\mathbf{p}}_-]) = 0.
\]
The operator of the spin magnetic moment
\begin{equation}
\hat{\boldsymbol{\mu}}_s = \mu_0(\boldsymbol{\sigma}_+ - \boldsymbol{\sigma}_-)
\tag{84.3}
\end{equation}
is not proportional to the operator of total spin $\hat{\mathbf{S}} = \frac{1}{2}(\boldsymbol{\sigma}_+ - \boldsymbol{\sigma}_-)$, and the operators $\hat{\mathbf{S}}^2$ and $\hat{\boldsymbol{\mu}}^2$ are not commutative. Therefore the states with definite values of the total spin $S$ and its projection $S_z$ are not eigenstates of the magnetic moment.

The states with given $S$ and $S_z$ are described by the spin functions $\chi_{SS_z}$ of the form
\begin{equation}
\begin{split}
\chi_{11} &= \alpha_+\alpha_-,\qquad\chi_{1,-1} = \beta_+\beta_-,\\
\chi_{10} &= \frac{1}{\sqrt{2}}(\alpha_+\beta_- + \alpha_-\beta_+),\\
\chi_{00} &= \frac{1}{\sqrt{2}}(\alpha_+\beta_- - \alpha_-\beta_+),
\end{split}
\tag{84.4}
\end{equation}
where $\alpha$ and $\beta$ are the spin functions of one particle, corresponding to the spin projections $+1/2$ and $-1/2$ (the indices ``$+$'' or ``$-$'' indicate whether the function refers to the positron or the electron). Of these the first two ($\chi_{11}$ and $\chi_{1,-1}$) are also eigenfunctions of the operator $\mu_z$, with eigenvalue $\mu_z = 0$. The func-

\SourcePage{386}{391}
tions $\chi_{10}$ and $\chi_{00}$, however, are not eigenfunctions of $\mu_z$; such are the combinations
\[
\frac{1}{\sqrt{2}}(\chi_{10}+\chi_{00}) = \alpha_+\beta_-,\qquad \frac{1}{\sqrt{2}}(\chi_{10}-\chi_{00}) = \alpha_-\beta_+.
\]

It is easy to see that the only non-zero matrix elements of the matrix $\langle S'S_z'|\mu_z|SS_z\rangle$, computed with the functions~(84.4), are the elements
\begin{equation}
\langle 00|\mu_z|10\rangle = \langle 10|\mu_z|00\rangle = 2\mu_0.
\tag{84.6}
\end{equation}

In weak magnetic fields (when $\mu_0 H \ll \Delta$, where $\Delta$ is the difference between the energies of levels with $S = 0$ and $S = l$) the initial approximation for computing the Zeeman splitting are the states with definite values of the total spin. In the first approximation this splitting is given by the average value of the energy of perturbation
\[
\hat{V}_H = -\hat{\mu}_z H.
\]
But all diagonal matrix elements of the operator $\hat{\mu}_z$, computed with the functions~(84.4), are zero. Thus, in weak fields the linear Zeeman effect in positronium is absent.

In the limiting case of strong fields ($\mu_0 H \gg \Delta$) one can neglect the spin interaction, leading to the establishment of definite values of~$S$. The components of the split level will in this case correspond to states with definite values of $\mu_z = \pm 2\mu_0$ (described by the functions~(84.5)), and the magnitude of their shift is $\pm 2\mu_0 H$.

\textbf{Problems}

\textbf{1.} Determine the fine structure of the levels of parapositronium (\textit{B.\,B.\,Berestetskii}, 1949)\,\footnote{The computation of the fine structure of orthopositronium see \textit{Sokolov A.\,A., Tsitovich V.\,N.}\/\/JETP.\ 1953.\ V.\,24.\ P.\,253.}.

\textbf{Solution.} The desired energy of level splitting is given by the average values of the correction terms in the Hamiltonian~(84.1), computed with the wave functions of the unperturbed states with $j = l$ (for $S = 0$ the contribution that differs from zero arises only from $\hat{V}_1$ and the second member in~$\hat{V}_3$).

The unperturbed wave functions (we denote them by $\psi$) satisfy the Schr\"odinger equation
\[
\hat{\mathbf{p}}^2\psi = -\Delta\psi = \left(E+\frac{1}{r}\right)\psi,\qquad E = -\frac{1}{4n^2}.
\]
Therefore
\[
\hat{\mathbf{p}}^4\psi = \hat{\mathbf{p}}^2\left(E+\frac{1}{r}\right)\psi = \left(E+\frac{1}{r}\right)^2\psi - \psi\Delta\frac{1}{r}-2\left(\boldsymbol{\nabla}\frac{1}{r}\right)(\boldsymbol{\nabla}\psi) =
\]
\[
= \left(E+\frac{1}{r}\right)^2\psi + 4\pi\delta(\mathbf{r})\psi + \frac{2}{r^2}\frac{\partial\psi}{\partial r}
\]

\SourcePage{387}{392}
The average value
\[
\overline{\mathbf{p}^4} = \overline{\left(E+\frac{1}{r}\right)^2}+4\pi|\psi(0)|^2 + \int_0^\infty\int\frac{\partial|\psi|^2}{\partial r}dr\,do.
\]
The last integral is non-zero only when $\psi(0) \neq 0$, i.e.\ only for $l = 0$, and the wave functions of $S$-states are spherically symmetric; the integral equals $-4\pi|\psi(0)|^2$ and cancels with the second term.

Introducing the operator of orbital angular momentum $\hat{\mathbf{l}} = [\mathbf{r}\hat{\mathbf{p}}]$ and writing
\[
-\hat{\mathbf{p}}^2 = \frac{\partial^2}{\partial r^2}+\frac{2}{r}\frac{\partial}{\partial r}-\frac{\hat{\mathbf{l}}^2}{r^2},\qquad -\hat{\mathbf{p}}^2\psi = -\left(E+\frac{1}{r}\right)\psi.
\]
From this, for the other mean value we need:
\[
\int\psi^*\frac{r}{r^3}(\mathbf{r}\hat{\mathbf{p}})\hat{\mathbf{p}}\psi\,d^3x = -\int\psi^*\frac{1}{r}\frac{\partial^2\psi}{\partial r^2}\,d^3x =
\]
\[
= \frac{1}{r}\left(E+\frac{1}{r}\right) - 4\pi|\psi(0)|^2 - l(l+1)\overline{r^{-3}}.
\]

(for $l = 0$ the last member is absent).

According to the well-known formulas of the theory of the hydrogen atom (see~III,~(36.14), (36.16)) with the replacement of the electron mass $m$ by $m/2$ we have
\[
|\psi(0)|^2 = \frac{1}{8\pi n^3}\delta_{l0},\qquad \overline{r^{-2}} = \frac{1}{2n^3(2l+1)},\qquad \overline{r^{-3}} = \frac{1}{4n^3 l(l+1)(2l+1)}\quad(l \neq 0).
\]

With the help of these formulas we finally obtain for the desired energy levels of parapositronium
\[
E_{nl} = -\frac{1}{4n^2}-\alpha^2\frac{me^4}{\hbar^2}\,\frac{1}{2n^3}\left(\frac{1}{2l+1}-\frac{11}{32n}\right).
\]

\textbf{2.} Determine the difference of energies of the ground states ($n = 1$, $l = 0$) of ortho- and parapositronium.

\textbf{Solution.} The dependence of energy on the total spin $S$ at $l = 0$ is contained only in the average value of the second member in $\hat{V}_3$ (the first member vanishes on averaging in a spherically-symmetric $S$-state)\,\footnote{The averaging over angles should be carried out before integration over $r$, as is seen from the manner of computation of the integral~(83.14), leading to the first member in~$\hat{V}_3$.}. The ground level of orthopositronium (${}^3S_1$) lies above the ground level of parapositronium (${}^1S_0$) by the amount
\[
E({}^3S_1)-E({}^1S_0) = \frac{7}{12}\alpha^2\frac{me^4}{\hbar^2} = 8.2\cdot 10^{-4}\,\text{eV}.
\]
